\section{Introduction}
\label{sec:introduction}

The minimum weighted feedback arc set (MWFAS) problem asks for a minimum-weight set of arcs whose removal makes a directed graph acyclic; equivalently, it asks for a vertex ordering with minimum total backward-arc weight. The problem is NP-hard \cite{K72} and arises whenever directed pairwise relations contain inconsistent cycles, including precedence aggregation, dependency repair, scheduling, and ranking from directed evidence \cite{F90,ACN08}. In weighted settings, not all violations are interchangeable: reversing a high-weight precedence relation is more costly than reversing a weak one, so useful heuristics must react to both graph structure and arc-weight asymmetry.

This paper focuses on \emph{sparse nonnegative weighted digraphs}. That regime differs structurally from dense linear-ordering or tournament instances. In many sparse graphs, large regions are already acyclic or only weakly cyclic, while the hard part of the objective is concentrated in a smaller set of strongly connected components (SCCs). Dense complete ordering instances behave differently: every ordered pair carries weight, the cyclic core effectively spans the whole graph, and matrix-based pairwise-ordering heuristics are naturally aligned with the input model. A method that is strong on sparse digraphs therefore need not be strongest on dense ordering benchmarks.

This creates a practical gap. Exact formulations provide certification but do not scale broadly enough to serve as the main tool on the sparse benchmark studied here \cite{BSNA21}. Classical greedy heuristics remain fast and interpretable, yet they may spend effort too globally or may not exploit the fact that residual backward weight is often localized after a strong constructive pass. Earlier local-ratio and weighted minimal/stable constructions provide useful seeds, but they leave open whether a \emph{protected} strongly connected component local-refinement stage can improve solutions consistently without ever returning a worse answer than the best available seed.

The central contribution is an \emph{incumbent-protected, component-local heuristic}. It begins from the better of two attributed seeds, scores nontrivial strongly connected components by their current backward-weight contribution, and applies component-local destroy-and-repair moves inside fixed original-graph neighborhoods. The destroy phase reactivates heavy removed internal arcs and removes light active internal arcs using deterministic prefixes; the repair phase reruns local-ratio reduction and heavy-first add-back only inside the selected component neighborhood. A candidate is accepted only if the repaired global ordering is feasible and strictly improves backward weight; otherwise the incumbent is restored exactly. We refer to this heuristic as incumbent-protected SCC neighborhood search (IPSNS). The new contribution is therefore not a new seed, but a non-worsening refinement layer for sparse digraphs: it performs localized destroy-and-repair inside selected cyclic components and can strengthen a competitive incumbent without ever returning a worse final solution.

The two seeds are supporting components rather than coequal contributions. LR-TA (local-ratio topological add-back) is an engineered sparse-graph seed derived from the directed local-ratio lineage of Demetrescu and Finocchi \cite{DF03} and connected to earlier author work on ranking-oriented MWFAS \cite{VahidiKoutis2024arxiv}. The weighted minimal/stable-style (WMSF-style) seed is an adapted implementation of the weighted minimal/stable feedback arc set line of Cavallaro and collaborators \cite{CCP23CEUR,CCP24,CC25}. The present paper does not claim those constructive foundations as new; it uses them as complementary entry points whose better incumbent IPSNS then protects and selectively improves.

The empirical message is intentionally narrow. On the primary sparse comparison, IPSNS attains the minimum observed backward weight among the evaluated methods on 96 of 97 standard nonnegative instances, with ties credited to every tied method (14 strict wins). A paired 97-instance comparison against the better seed yields 14 strict improvements, 83 ties, and 0 regressions. On the 57-instance exact-validation subset, the mean optimum-normalized gap is 0.0031\%. On the 93-instance common sparse subset, repeated-run per-instance medians over 20 runs yield 38 wins, 55 ties, and 0 losses for IPSNS against the evaluated DRMacIver/FAS executable. A single-run dense LOLIB transfer study shows that the selected complete-ordering benchmark and evaluated implementations favored DRMacIver/FAS on 45 of 50 instances. The paper's claim is therefore that protected SCC-local refinement is an effective sparse-graph algorithm-engineering strategy among the evaluated methods, not that it solves every MWFAS regime equally well.

The contributions are:
\begin{itemize}
    \item \textbf{Problem focus.} The paper studies minimum weighted feedback arc set on sparse nonnegative weighted digraphs and keeps that regime separate from dense linear-ordering and tournament-style benchmarks.
    \item \textbf{Algorithmic contribution.} We introduce IPSNS, an incumbent-protected strongly connected component local-refinement layer. Starting from the better of two attributed seeds, it applies localized destroy-and-repair inside selected cyclic component neighborhoods, reruns local-ratio reduction and heavy-first add-back there, and accepts only strict global improvements. Therefore the method cannot return a worse solution than its incumbent.
    \item \textbf{Attribution and supporting analysis.} The constructive seeds are explicitly attributed to prior local-ratio and weighted minimal/stable ideas rather than claimed as new. We formalize implementation-level properties for seed feasibility, one-pass add-back minimality, and IPSNS incumbent non-worsening, together with concise complexity summaries.
    \item \textbf{Computational contribution.} We report a reproducible sparse-digraph evaluation including seed comparisons, transparent weighted calibrators, a graph-library baseline, an external executable comparator, exact and MIP validation, repeated-run summaries, and a dense LOLIB transfer study used only as a scope boundary.
\end{itemize}

The value of the refinement is not that every instance improves; rather, it strengthens already competitive constructive orderings without regressions, identifies the remaining improvable sparse instances, and gives a reproducible account of where component-local search helps and where dense ordering methods remain preferable.

The remainder of the paper positions the work, defines the objective, presents the framework, summarizes the experimental design, reports the results, discusses the sparse and dense regimes, and concludes.
